\documentclass[../main.tex]{subfiles}
\begin{document}

\chapter*{Exercises Lecture 1}

\section{Sampling random points within D-dimensional domains by hit and miss}
Use the hit-and-miss method to estimate numerically the volume of an ellipsoid of parameters $a=3$, $b=2$, $c=2$ and plot the deviation of the estimate from the analytic value as a function of the number of throws used.
Repeat the above procedure to estimate the volume of an ellipsoid of parameters $a=3$, $b=1$, $c=1$. Compare the two results and try to improve the second case if needed.

\subsection{Solution}
By randomly sampling points in the volume [-a, a] x [-b, b] x [-c, c] we get the figure \ref{fig: ex_1_1}. Each point in the figure is calculated as $|V_{exact} - V_{estimated}|$ where $V_{estimated}$ is the average of 25 runs of hit and miss.
At $10^7$ points we get:
\begin{verbatim}
Exact volume for ellipsoid (3,2,2): 50.2655, best estimate: 50.26699
Exact volume for ellipsoid (3,1,1): 12.5664, best estimate: 12.56602
\end{verbatim}


\begin{figure}[ht!]
    \centerline{\includesvg[width=0.95\columnwidth]{/imgs/ex_1_1.svg}}
    \caption{Sampling results for the estimate of the volume of an ellipsoid.}
    \label{fig: ex_1_1}
\end{figure}


%%%%%%%%%%%%%%%%%%%%%%%%
% \begin{wrapfigure}{r}{0.5\textwidth}
%   \centering
%   \includesvg[width=0.33\textwidth]{./imgs/ex_1_1.svg}
% \caption{Explanation of the content ..............}
% \label{fig:1.1}
% \end{wrapfigure}
% %%%%%%%%%%%%%%%%%%%%%%%%%


\section{Sampling random numbers from a given distribution: inversion method}
Use the inversion method to generate random numbers with the following PDF:
\begin{enumerate}
    \item $\rho(x) = c x e^{-x^2}$, for $x \in \mathbb{R}^+$.
    \item $\rho(x) = b x^4$, for $x \in [0, 3]$.
\end{enumerate}

\subsection{Solution}
First we find the normalization constant by imposing $\int \rho(x)dx = 1$, so we obtain 
\begin{align}
\rho_1(x) = 2 x e^{-x^2} && \rho_2(x) = \frac{5}{243}x^4
\end{align}
then we calculate the cumulative $F(x)$:
\begin{align}
F_1(x) = \int_0^x 2te^{-t^2}dt = 1-e^{-x^2} && F_2(x) = \int_0^3\frac{5}{243}t^4dt = \frac{x^5}{243}
\end{align}
and finally we obtain $F^{-1}(y)$ by inverting $F(x)$:
\begin{align}
y = 1-e^{-x^2} \implies x = F^{-1}_1(y) = \sqrt{\log{\frac{1}{1-y}}}\ , \text{for } y \in [0, 1) \\
y = \frac{x^5}{243} \implies x = F^{-1}_2(y) = \sqrt[5]{243y}\ , \text{for } y \in [0, 1]
\end{align}
By sampling uniformly $y$ and then calculating $F^{-1}(y)$ it's like sampling $\rho(x)$:

\begin{figure}[ht!]
    \centerline{\includesvg[width=0.99\columnwidth]{/imgs/ex_1_2.svg}}
    \caption{Comparison of the sampling with the real pdfs. On the left $\rho_1$, on the right $\rho_2$} 
    \label{fig: ex_1_2}
\end{figure}
\end{document}


